% ----------------------------------------------------------
% RESUMOS
% ----------------------------------------------------------
\begin{resumo}
O presente projeto tem como objetivo desenvolver um sistema web para o gerenciamento e controle de vendas e estoque para a loja \textbf{Tropical Mix}, localizada em Guarulhos. A loja em questão realiza vendas de insumos para máquinas de sorvete expresso e fabricação de sorvetes. O sistema visa otimizar a gestão dos produtos, atualizar e monitorar o inventário e manter um registro das transações comerciais realizadas na loja. 

O projeto foi idealizado para atender às necessidades de gestão da loja que enfrentam dificuldades no acompanhamento manual de entradas e saídas de produtos. A solução utiliza tecnologias como \textbf{Python (Flask)}, \textbf{MySQL} e \textbf{React}, promovendo uma interface simples, intuitiva e acessível por navegadores web. Além disso, o sistema provê relatórios de análise de dados com insights a fim de facilitar as tomadas de decisão.

A iniciativa possibilita aos estudantes unir conhecimentos teóricos adquiridos ao longo da graduação, aliando teoria e prático no desenvolvimento do sistema e processos gerenciais do mesmo.

\vspace{\onelineskip}

\noindent
\textbf{Palavras-chave}: Sorveteria. Sistema de Gerenciamento de Estoque. Histórico de vendas.
\end{resumo}

\begin{resumo}[Abstract]
\begin{otherlanguage*}{english}
This project aims to develop a web system for managing and controlling sales and inventory for the store \textbf{Tropical Mix}, located in Guarulhos. The store sells supplies for soft-serve ice cream machines and ice cream production. The system seeks to optimize product management, update and monitor inventory, and maintain records of commercial transactions carried out in the store.

The project was designed to meet the management needs of the store, which faces difficulties in manually tracking product inflows and outflows. The proposed solution employs technologies such as \textbf{Python (Flask)}, \textbf{MySQL}, and \textbf{React}, providing a simple, intuitive interface accessible through web browsers. Furthermore, the system offers data analysis reports and insights to support decision-making processes.

This initiative allows students to integrate theoretical knowledge acquired during their studies with practical experience, combining theory and practice in the development of the system and its management processes.

\vspace{\onelineskip}

\noindent
\textbf{Keywords}: Ice Cream Shop. Inventory Management System. Sales History.
\end{otherlanguage*}
\end{resumo}
